\section{Conceptos basicos}

\subsection{Que hace el RCE}

El RCE permite escribir reglas CQL dentro de la misma pantalla del sistema,
validarlas contra CQL Translation Service, guardarlas como artefactos FHIR en
HAPI y evaluarlas contra pacientes sinteticos. Cuando una regla aplica, el
usuario ve una card CDS educativa.

\subsection{Modo aula anonimo}

No hay login visible. Cuando un navegador entra al RCE, el backend crea o
recupera una sesion anonima mediante una cookie segura. Esa sesion apunta a un
sandbox.

Cada sandbox separa:

\begin{itemize}
  \item reglas creadas por el alumno;
  \item cambios hechos sobre pacientes;
  \item cards generadas;
  \item historial de actividad CDS.
\end{itemize}

Asi varios alumnos pueden abrir el mismo paciente base sin pisarse entre ellos.
El paciente sintetico base vive en HAPI; los cambios del alumno quedan como
overlay de su sandbox.

\subsection{Roles Alumno y Docente}

El modo Alumno sirve para practicar: navegar pacientes, escribir reglas, probar
reglas y ver cards dentro de su sandbox.

El modo Docente habilita acciones de clase como publicar, activar o desactivar
reglas compartidas. El cambio de rol es parte de la maqueta educativa; no es un
sistema de autenticacion institucional.

\rceScreenshot{04-sandbox-sesion.png}{Menu de sesion anonima con aula, sandbox y rol actual.}

\subsection{Datos clinicos}

Los datos se representan con recursos FHIR R4. HAPI FHIR es el repositorio
logico y PostgreSQL es solo la persistencia interna de HAPI. El backend no entra
a las tablas internas de PostgreSQL.

Los pacientes de la demo son sinteticos y pueden incluir recursos como
\texttt{Patient}, \texttt{Condition}, \texttt{Observation}, \texttt{MedicationRequest}
y \texttt{Encounter}.

\subsection{Por que la ficha tiene condiciones, observaciones, medicamentos y encuentros}

Las pestanas de la ficha no son categorias inventadas solo para la maqueta.
Representan recursos del estandar HL7 FHIR R4. Esto es importante para la clase:
una regla CQL real no deberia leer campos sueltos de una pantalla, sino consultar
recursos clinicos interoperables.

\begin{longtable}{p{0.20\linewidth}p{0.28\linewidth}p{0.42\linewidth}}
\toprule
Pestana en el RCE & Recurso FHIR & Para que sirve en reglas CQL \\
\midrule
Resumen & \texttt{Patient} + resumen calculado & Identifica al paciente, edad, sexo, fecha de nacimiento y conteos utiles para orientar la ficha. \\
Condiciones & \texttt{Condition} & Diagnosticos, problemas o condiciones activas. Ejemplo: diabetes, hipertension, asma. \\
Observaciones & \texttt{Observation} & Laboratorios, signos vitales y mediciones. Ejemplo: HbA1c, presion arterial, peso, glicemia. \\
Medicamentos & \texttt{MedicationRequest} & Ordenes o indicaciones de medicamentos. Ejemplo: metformina activa o ausencia de tratamiento. \\
Encuentros & \texttt{Encounter} & Consultas, controles, visitas, hospitalizaciones o contexto de atencion. \\
CDS & \texttt{CDS Hooks Cards} & Recomendaciones generadas al evaluar reglas contra el paciente. \\
\bottomrule
\end{longtable}

\subsection{Son necesarios para todas las reglas}

No. Depende de lo que pregunte la regla.

Una regla por edad usa principalmente \texttt{Patient}:

\begin{lstlisting}
context Patient

define "Aplica":
  AgeInYears() >= 18
\end{lstlisting}

Una regla por laboratorio necesita \texttt{Observation}. Por ejemplo, HbA1c:

\begin{lstlisting}
[Observation: "Hemoglobin A1c"] O
  where O.status.value = 'final'
    and (O.value as FHIR.Quantity).value.value >= 6.5
\end{lstlisting}

Una regla por diagnostico necesita \texttt{Condition}. Por ejemplo, buscar si el
paciente tiene una condicion activa compatible con diabetes.

\begin{lstlisting}
[Condition] C
  where C.clinicalStatus.coding[0].code.value = 'active'
\end{lstlisting}

Una regla de medicamentos puede combinar \texttt{Condition} y
\texttt{MedicationRequest}. Ejemplo docente: paciente con diabetes que no tiene
una indicacion de metformina.

Una regla relacionada con el momento de atencion puede usar \texttt{Encounter}.
Ejemplo docente: mostrar una card solo durante una consulta ambulatoria o en un
control reciente.

\subsection{Como explicarlo a los alumnos}

La idea pedagogica es esta:

\begin{lstlisting}
FHIR guarda los hechos clinicos
CQL expresa la logica de la regla
CDS Hooks define cuando se pide ayuda
La card muestra la recomendacion
\end{lstlisting}

Por eso el RCE muestra esos grupos clinicos. No estan ahi por decoracion:
permiten ensenar que una regla puede activarse por edad, por diagnosticos, por
laboratorio, por signos vitales, por medicamentos o por el contexto de una
atencion.

\subsection{Identificadores tecnicos}

FHIR necesita identificadores estables para leer y relacionar recursos, por
ejemplo \texttt{Patient/123}, \texttt{Observation/abc} o un
\texttt{correlationId}. En el RCE esos IDs se muestran de forma compacta para no
ensuciar la ficha. El nombre, edad, sexo, fecha de nacimiento y resumen clinico
son la informacion principal para el usuario; el ID queda como metadata tecnica
para depuracion, trazabilidad y pruebas por terminal.

\subsection{Referencias oficiales}

\begin{itemize}
  \item \url{https://hl7.org/fhir/R4/patient.html}
  \item \url{https://hl7.org/fhir/R4/condition.html}
  \item \url{https://hl7.org/fhir/R4/observation.html}
  \item \url{https://hl7.org/fhir/R4/medicationrequest.html}
  \item \url{https://hl7.org/fhir/R4/encounter.html}
\end{itemize}
